% ============================================================================
%  12_conclusiones.tex — Conclusiones.
% ============================================================================
\section{Conclusiones}
\label{sec:conclusiones}

Se han cumplido los objetivos de la Actividad~3 y se ha cerrado el ciclo de
desarrollo del producto inteligente de \emph{ACME\,S.A.}. La base fusionada
Act1+Act2 incorpora ahora una capa de \textbf{instrumentación avanzada} que
aporta cuatro funciones de autodiagnóstico ---redundancia con conmutación
automática, detección de averías por límites de rango, promediado de medidas
y vigilancia de la corriente de actuadores--- y reacciona ante el fallo con
una jerarquía escalonada \textsf{OK}/\textsf{AVISO}/\textsf{ALARMA} que
distingue entre situaciones con respaldo y fallos críticos. La diagnosis de
temperatura y humedad es \emph{independiente}, reflejando que el DHT22 tiene
elementos de medida distintos para cada magnitud.

La solución encaja con los cuatro criterios de evaluación: \emph{(C1)}\
se han probado y documentado las partes del ejemplo de autodiagnóstico, con
trazas CSV reales del sistema en sus distintos modos de fallo;
\emph{(C2)}\ el circuito E/S se ha ampliado con dos componentes nuevos ---un
DHT22 de reserva y un sensor de corriente acondicionado--- conectados a los
únicos pines que quedaban libres en el ESP32 tras Act1+Act2;
\emph{(C3)}\ el funcionamiento es eficiente y la programación mínima:
\emph{toda} la capa nueva vive en un único módulo \texttt{diagnostics} y el
resto de módulos heredados se han tocado lo justo (sólo un \emph{getter}
adicional en \texttt{elevator} para forzar el estado \textsf{ALARM});
\emph{(C4)}\ la presentación sigue el formato pedido (Carlito ≡
Calibri 12, interlineado 1.5) y respeta las diez páginas.

Como nota explícita: la \textbf{planta~4} se ha quedado sin pulsador físico
para que el GPIO\,17 acoja el segundo DHT22, pero sigue siendo accesible por
el botón~4 del mando IR ---pulsadores y mando se diseñaron redundantes desde
Act2---, por lo que la operativa no se degrada.

\paragraph{Mejoras futuras.} Sustituir el potenciómetro por un módulo INA219
real manteniendo la misma API de \texttt{diagnostics}; añadir un
\emph{watchdog} hardware; y exportar la telemetría por Wi-Fi a un servicio
externo, dando el salto a una instrumentación conectada propia del IoT
industrial.
